% !TeX root = main.tex
\chapter{Introduction\label{chap:1}}
\noindent \LaTeX (pronounced "LAY-tek" or "LAH-tek") is a tool for typesetting professional-looking documents. However, LaTeX's mode of operation is quite different to many other document-production applications you may have used, such as Microsoft Word or LibreOffice Writer.

\section{Why learn \LaTeX}

\noindent Various arguments can be proposed for, or against, learning to use LaTeX instead of other document-authoring applications; but, ultimately, it is a personal choice based on preferences, affinities, and documentation requirements. \\

Arguments include:
\begin{enumerate}
\item[1.] Support for typesetting extremely complex mathematics, tables and technical content for the physical sciences;
\item[2.] Facilities for footnotes, cross-referencing and management of bibliographies;
\item[3.] Ease of producing complicated, or tedious, document elements such as indexes, glossaries, table of contents, lists of figures;
\item[4.] Being highly customizable for bespoke document production due to its intrinsic programmability and extensibility through thousands of free add-ons.
\end{enumerate}

\noindent For more, check:

\href{https://www.overleaf.com/learn/latex/Learn_LaTeX_in_30_minutes#What_is_LaTeX?}{Overleaf Tutorial}




\section{Another Part of the Introduction}

\noindent \LaTeX (pronounced "LAY-tek" or "LAH-tek") is a tool for typesetting professional-looking documents. However, LaTeX's mode of operation is quite different to many other document-production applications you may have used, such as Microsoft Word or LibreOffice Writer.\\

\noindent \LaTeX (pronounced "LAY-tek" or "LAH-tek") is a tool for typesetting professional-looking documents. However, LaTeX's mode of operation is quite different to many other document-production applications you may have used, such as Microsoft Word or LibreOffice Writer.\\

\section{Problem Formulation}\label{sec:problem}
\noindent It's probably a good time to start focusing in on the issue.\\

\section{Knowledge Gaps}\label{sec:gaps}
\noindent What do we \emph{not} know:
\begin{itemize}
	\item This issue with regards to:
	\begin{itemize}
		\item Particle physics
		\item Dark matter
	\end{itemize}
	\item With emphasis on methods for determining the exact velocity of a water drop.
\end{itemize}